% Caption for the architecture figure.
\caption{%
  Proposed degradation-aware, amplitude-focused bandwidth extension with a Schr\"odinger
  bridge. This is a single architecture, not an ablation: the trunk is a selective state
  space model in both streams, and every component shown is active in the trained model.
  Blue blocks are cited: STFT/iSTFT, the anti-wrapping phase losses and the MRAD/MRPD
  discriminators follow AP-BWE~\cite{lu2024apbwe}; the Lattice coupling follows
  CIS-BWE~\cite{tamiti2026cisbwe}. Orange blocks with a numbered badge are this work.
  \badgenum{1}~a blind encoder infers an eight-dimensional degradation descriptor from
  the narrowband input alone; its embedding $\mathbf e$ conditions every Lattice stage.
  \badgenum{2}~its cutoff estimate $k_c$ builds a coarse wideband (SBR) prior that serves
  three roles at once: a second input channel, the base the magnitude head predicts a
  residual on, and the bridge endpoint $x_0$.
  adgenum{3}~a selective state space trunk in both streams, with separate weights per stream.
  \badgenum{4}~quantum-inspired amplitude modules: a state-encoding angle/amplitude
  embedding at the input and a tensor-train (MPS) output head that is
  $3.8\times$ smaller than the dense projection it replaces.
  \badgenum{5}~a Schr\"odinger bridge between the observed and clean pairs, with phase
  bridged on the unit circle. Because the SBR prior is the endpoint, the bridge corrects
  a band that is already filled rather than inventing one. Under the data-prediction
  parameterisation $N{=}1$ reduces exactly to the feed-forward model, so the iterative
  cost is opt-in.
  \badgenum{6}~the phase path is driven by instantaneous frequency and group delay and
  refined by a 2D, frequency-dilated head on baseband phase differences, which removes
  the deterministic $2\pi mH/N$ ramp~\cite{gaznepoglu2026vocos}; it is held at a setting
  whose APD already matches AP-BWE's.
  \badgenum{7}~an amplitude-heavy critic set: UHD scores the log-frequency magnitude
  envelope and SBMD returns per-band local verdicts, with BiCoD as a cheap phase floor
  that asks whether the synthesised band is phase-locked to the observed one.
  \badgenum{8}~an output-side baseband graft. Everything below $k_c$ was handed to the
  model, so it is restored through a raised-cosine crossover rather than re-synthesised.
  \badgenum{9}~three amplitude losses that the usual log-amplitude MSE cannot supply:
  spectral convergence is relative and therefore scale-sensitive, so a uniform gain error
  cannot be carried for free; per-band energy in dB is the coarse-grained form of the LSD
  metric itself; and a multi-resolution log-magnitude $L_1$ constrains spectral shape.
  On 4\,kHz$\to$16\,kHz VCTK after
  6{,}250 steps on 2{,}000 training files it reaches LSD 1.20, NISQA-MOS 3.39, PESQ 2.20, STOI 0.925 and
  SI-SDR 11.7\,dB on the official test speakers, with checkpoint selection on held-out training speakers.%
}
